% ==============================================================================
%  Chapter 7 -- Glueball spectroscopy with a learned variational operator
%  Purpose: the headline-result chapter, ~25-30 pages. This is DONE physics:
%  nearly all of it can be lifted (with thesis-level expansion) from
%  glueball_report/glueball_spectroscopy.tex and
%  notes/glueball_spectroscopy.md. Numbers below are final unless a third
%  ensemble (seed 2) is added.
% ==============================================================================

\section{The idea: a network as a variational operator}
\label{sec:glue-idea}

\TODO{Frame shift from Ch.~6: instead of regressing a known observable,
train GELT to BE an operator. \texttt{GELT(reduction="none")} emits a
per-site scalar field $\to$ zero-momentum projection $\to$ timeslice
operator $\Obar(t)$. The Rayleigh-quotient loss $-C(1)/C(0)$ is maximised
exactly when the operator's ground-state overlap $A_0$ is maximised, and
at the optimum the converged loss IS $e^{-m \at}$: the loss value is a
mass measurement. This is Morningstar--Peardon's variational principle
\cite{morningstar1999} with a neural network replacing the operator
basis.}

\section{The classical baseline}
\label{sec:glue-classical}

\subsection{Toolkit}
\TODO{Spatial-only APE smearing (gauge covariant, unit-tested; time links
untouched so the transfer-matrix reading holds), the $0^{++}$ operator
(sum of spatial-plane Wilson loops), vacuum-subtracted correlator,
multi-level GEVP over cumulative smearing levels (eigh-whitening with an
eigenvalue floor -- low statistics can make $C(t_0)$ indefinite),
jackknife everywhere.}

\subsection{Run-by-run: the isotropic dead end and the anisotropic fix}
\TODO{Tell it honestly as a negative-then-positive result. Isotropic
$L=12$, $\beta=2.4$, $N=2000$: no plateau -- weak $m a \approx 0.8$
drowns by $\Delta \approx 3$ even with the GEVP; the LATTICE, not the
operator basis, was the bottleneck. Anisotropic $L=12$, $L_t=24$,
$\beta=2.4$, $\xi=3.0$, $N=2000$: clean decay over 10--12 slices,
GEVP ground state plateaus at
$\meff(\Delta{=}1) = 0.365 \pm 0.008$,
$\meff(\Delta{=}2) = 0.333 \pm 0.011$, i.e.\
$m\,\at \approx 0.33$ -- THE ANCHOR every learned result is judged
against. Caveat box: $m\,\as = \xi\, m\,\at \approx 1.0$ at
$\beta_s = 0.8$ is strong-coupling, not continuum physics; no $\xi$
tuning, no continuum extrapolation (\cref{sec:concl-future}).}

\section{Constraints on the learned operator}
\label{sec:glue-constraint}

\TODO{The theory-driven design decisions (from the 2026-07-01 audit) --
worth a full section because they are where the physics disciplines the
ML: (i) the operator must be PER-TIMESLICE 3D (\texttt{GELT(D=3)} on each
timeslice's spatial links, batch = config $\times$ timeslice) -- any
temporal receptive field voids the variational bound and makes the
Rayleigh loss gameable toward $m \to 0$ ("never smear in time" applied to
the network itself); (ii) hard train/held-out config split, blocked
jackknife on held-out configs only; (iii) \texttt{mlp\_zero\_init} must
be OFF (zero init $\Rightarrow$ identically zero Rayleigh gradient).}

\subsection{Making the loss trainable}
\TODO{The two loss pathologies and their fixes (real lessons, keep them):
chained-ratio denominators replaced by $C(0)$; and the $\Obar \to
\lambda\Obar$ flat direction, which site-coherent smeared inputs pump to
float32 overflow, pinned by a $(\log C(0))^2$ scale regulariser.}

\section{Training runs}
\label{sec:glue-runs}

\subsection{Run 4 --- thin input: a genuine operator that loses}
\TODO{Thin-plaquette input: GELT learns $\approx$ APE$\times 2$'s staple
content (a depth-4 loop degree cannot rebuild iterated smearing) and is
variationally redundant with the classical basis. Valuable as a
capacity-vs-physics lesson.}

\subsection{Run 5 --- smeared input channels: the physics fix}
\label{sec:glue-run5}
\TODO{Multi-level smeared input channels
(\texttt{INPUT\_SMEAR\_LEVELS} $= (0,2,4,6)$,
$3 \times n_{\mathrm{levels}}$ input channels). Results: Rayleigh loss
saturates the transfer-matrix bound at the anchor
(val $-0.6185 \leftrightarrow m\,\at \approx 0.33$); $\meff$ plateaus
from $\Delta = 1$; beats the classical GEVP at $\Delta = 1$ by
$3.9\sigma$ (correlated blocked jackknife of the difference on shared
test configs: $-0.028 \pm 0.007$); the enlarged 5-operator GEVP
(classical basis + GELT) collapses onto pure GELT -- the learned operator
already spans the useful directions.}

\section{Quoting the result: cosh fits and the overlap \texorpdfstring{$A_0$}{A0}}
\label{sec:glue-overlap}

\TODO{The LGT-standard presentation: masses from a correlated cosh fit
(profiled-$A$ grid fit, refit inside every jackknife sample), operator
quality as the ground-state overlap fraction
$A_0 = A(1 + e^{-m N_t})/C(0)$ -- exactly the quantity the Rayleigh loss
maximises. Comparator: the PROJECTED classical GEVP operator (fixed
ground-state vector $v_0$), the optimal single operator in the classical
basis's span -- the apples-to-apples opponent.
Numbers (window $\Delta \in [2,7]$, 400 held-out test configs):}
\begin{center}
\begin{tabular}{lcc}
\hline
Operator & $m\,\at$ & $A_0$ \\
\hline
GELT                  & $0.332 \pm 0.027$ & $0.903 \pm 0.047$ \\
GEVP-projected        & $0.340 \pm 0.030$ & $0.837 \pm 0.056$ \\
\hline
\multicolumn{3}{l}{correlated $\Delta A_0 = +0.066 \pm 0.031$
  ($2.1\sigma$), $\Delta m$ consistent with $0$}\\
\hline
\end{tabular}
\end{center}
\TODO{Interpretation sentence: same mass, more ground-state weight -- the
network is a better operator, not different physics. Figure:
\texttt{glueball\_overlap.png} ($\meff$ + fit bands; $\rho(\Delta)$
overlap panel, flat at $A_0 \Leftrightarrow$ pure ground state).}

\section{Replication on an independent ensemble}
\label{sec:glue-replication}

\TODO{Fresh ensemble (new seed, sampled from scratch) + from-scratch
training. Results: GELT $A_0 = 1.013 \pm 0.062$ vs.\ GEVP-projected
$0.925 \pm 0.071$; correlated $\Delta A_0 = +0.089 \pm 0.030$
($2.9\sigma$); $\meff$ difference at $\Delta = 1$: $-0.038 \pm 0.008$
($4.5\sigma$). Combined with Run 5, the headline:
$\boldsymbol{\Delta A_0 = +0.078 \pm 0.022}$ ($3.6\sigma$).
Key methodological lesson to write up: the Rayleigh floor is
ENSEMBLE-SPECIFIC -- the replication ensemble reads $\sim 1\sigma$
heavier (GEVP flat at $\approx 0.39$; best val
$-0.5638 \leftrightarrow m \approx 0.395$, saturated to $\sim 0.004$ like
Run 5's $-0.6185 \leftrightarrow 0.33$), so LOSS VALUES do not replicate
but OPERATOR QUALITY does. \TODO{if the seed-2 ensemble is rerun (the
first attempt died on the Creutz rejection cap, since fixed), add it as a
third independent measurement -- a bonus, not a blocker.}}

\section{Open items in this chapter}
\label{sec:glue-open}

\TODO{Declared before the referee asks: (i) the matched-parameter L-CNN
baseline on the SAME per-timeslice variational task -- the last missing
column of the shootout; (ii) anisotropy tuning / continuum limit;
(iii) a third ensemble.}
